\documentclass[../main.tex]{subfiles}
\graphicspath{{\subfix{../images/}}}
\begin{document}


\begin{figure}[tb]
  \centering
  \subfigure[single]{
    \includegraphics[width=0.23\textwidth]{single.png}
  }
  \subfigure[double]{
    \includegraphics[width=0.23\textwidth]{double.png}
  }
  \subfigure[ridge]{
    \includegraphics[width=0.23\textwidth]{ridge.png}
  }
  \centering
  \subfigure[diagonal]{
    \includegraphics[width=0.23\textwidth]{diagonal.png}
  }
  \subfigure[corners]{
    \includegraphics[width=0.23\textwidth]{corners.png}
  }
  \subfigure[triangle]{
    \includegraphics[width=0.23\textwidth]{triangle.png}
  }
  \centering
  \subfigure[circle]{
    \includegraphics[width=0.23\textwidth]{circle.png}
  }
  \subfigure[random]{
    \includegraphics[width=0.23\textwidth]{random.png}
  }
%  \subfigure[t=9]{
%    \includegraphics[width=0.23\textwidth]{frame91.jpg}
%  }
  \caption{Generated patterns}
  \label{fig:1}
\end{figure}

\subsection{Previous Work Summary}

Our initial research established a framework for surface pattern recognition on synthetic three-dimensional floor point clouds using Shape Subspace methods combined with Difference Subspace (DS) techniques. The methodology involved segmenting synthetic floor surfaces (50×50 units, 10,000 points) into uniform patches, generating shape subspaces via Singular Value Decomposition (SVD), and computing difference subspaces between neighboring patches to detect local surface variations.

Key findings from our previous work demonstrated that:
\begin{itemize}
    \item Smaller patch sizes (3×3, 4×4) provided higher accuracy for fine-grained surface detection
    \item Four-neighbor configurations delivered better sensitivity to local variations
    \item Eight-neighbor approaches offered improved stability for complex patterns
    \item The method successfully detected various synthetic patterns (bumps, ridges, corners)
\end{itemize}

\subsection{Motivation and Objectives}

Although our initial approach proved effective on synthetic data, limitations became apparent: (1) boundary sensitivity at patch edges, (2) limited spatial resolution from non-overlapping patches, and (3) unvalidated performance on real-world noisy data. 

This progress report addresses these through the development of a sliding window methodology with overlapping analysis regions, the validation of real PLY files (61,914 points), and the establishment of data-dependent method selection guidelines. The sliding window approach provides continuous coverage without rigid boundaries, bridging the gap between synthetic validation and real-world application.

\end{document}